\subsubsection{Bloque con vínculo elástico bajo impulso instantáneo}

Un bloque de masa $m$ está conectado a un resorte lineal de constante $k$ y longitud natural $\ell_0$ como se muestra en la \textbf{Figura \ref{fig:bloque_impulso}}. El bloque se encuentra inicialmente en reposo y el resorte sin estirar cuando se le aplica un impulso instantáneo $\hat{P}$ (modelado mediante una delta de Dirac). Asumiendo que la gravedad actúa verticalmente hacia abajo, realice lo siguiente:

\begin{enumerate}
    \item Defina los marcos de referencia y sistemas de coordenadas adecuados para el movimiento unidimensional.
    \item Analice la cinemática del bloque justo después de la aplicación del impulso para determinar su velocidad inicial.
    \item Analice la cinética del sistema mediante balances de energía para determinar la máxima compresión del resorte.
\end{enumerate}

\begin{figure}[!ht]
    \centering
    \begin{minipage}{0.7\textwidth}
        \centering
        \begin{tikzpicture}[>=Latex, scale=1.3]
            % --- Suelo y Pared ---
            \fill[pattern=north east lines] (-0.7,-0.2) rectangle (4.5,0);
            \draw[thick] (-0.5,0) -- (4.5,0);
            \fill[pattern=north east lines] (-0.7,0) rectangle (-0.5,1.5);
            \draw[thick] (-0.5,0) -- (-0.5,1.5);

            % --- Resorte ---
            \draw[decoration={aspect=0.3, segment length=4pt, amplitude=4pt, coil}, decorate] 
                (-0.5, 0.5) -- (1.5, 0.5);
            \node[above] at (0.5, 0.7) {$K$};
            \node[above right, xshift=-3pt, yshift=1pt] at (-0.5, 0.5) {$Q$};
            \filldraw[] (-0.5,0.5) circle (1pt);

            % --- Bloque ---
            \draw[thick, fill=gris1!20] (1.5, 0) rectangle (2.7, 1);
            \node at (2.1, 0.5) {$m$};
            \filldraw[] (1.5,0.5) circle (1pt);

            % --- Cotas ---
            \draw[|<->|] (-0.5, -0.4) -- (1.5, -0.4) node[midway, fill=white] {$\ell_0$};
            \draw[|->, thick] (2.1, -0.4) -- (1.5, -0.4) node[midway, below] {$x$};

            % --- Impulso ---
            \draw[<- ,ultra thick] (2.7, 0.5) -- (3.7, 0.5) node[right] {$\mathbf{\hat{P}}$};

            % --- Gravedad ---
            \draw[->, thick, myred] (4.2, 1.2) -- (4.2, 0.4) node[midway, right] {$\mathbf{g}$};
        \end{tikzpicture}
    \end{minipage}
    \caption{Esquema de un bloque de masa $m$ sujeto a un resorte bajo la acción de un impulso $\hat{P}$.}
    \label{fig:bloque_impulso}
\end{figure}